\documentclass[twocolumn]{aastex631}
\begin{document}
\title{The reionization ionizing-photon-budget shortfall for star-forming galaxies is not robust to systematics at z\~{}9 (optimistic systematic corner)}
\author{NebulaMind Lab (autonomous pipeline)}
\affiliation{NebulaMind Lab, \url{https://lab.nebulamind.net}}
\begin{abstract}
Reionization ionizing-photon-budget reconciliation at z\~{}9: star-forming galaxies require f\_esc=0.072 (+0.062/-0.033) to close the budget (Madau-Dickinson SFRD, log xi\_ion=25.5+/-0.15, clumping C=2-5, JWST-SFRD tail), versus indirect-proxy-inferred f\_esc=0.080 (+0.147/-0.051) from LzLCS O32/beta calibrations. Median delta(required-inferred)=-0.007 dex-frac (16-84\%: -0.150 to +0.067); 46\% of the systematic MC shows a shortfall. Conclusion: the budget CLOSES within the systematic, and the sign flips sign between the O32 and beta calibrations (calibration-driven, not robust). The result is bounded by the xi\_ion x clumping x proxy-calibration systematic, not by statistics. Generated autonomously from public data (jwst) via the NebulaMind Lab runner: a bounded, reproducible, descriptive study.
\end{abstract}
\section{Introduction}
The reionization-photon-budget crisis has sparked significant interest in understanding the role of star-forming galaxies during this period [Muoz2024]. Previous studies have highlighted the importance of accurately calibrating models to account for ionizing photons [Park2022, Davies2021]. In particular, reconciling the cosmic star formation rate density (SFRD) with the required ionizing photon budget is crucial for a comprehensive understanding of reionization. This study aims to address this challenge by examining the ionizing-photon-budget reconciliation at z\~{}9.
\section{Data and method}
To investigate this issue, we employ a literature-anchored budget calculation that does not rely on survey catalog data. Instead, we use the Madau \& Dickinson (2014) analytic fitting function for the cosmic SFRD and adopt published values for xi\_ion and O32/beta f\_esc proxy calibrations from LzLCS [Chisholm+22, Flury+22; Simmonds+24]. By focusing on systematics reconciliation over existing literature values, we aim to provide a clearer understanding of the reionization photon budget.
\section{Result}
Our analysis reveals that star-forming galaxies require an escape fraction f\_esc = 0.072 (+0.062/-0.033) to close the ionizing-photon-budget at z\~{}9, considering the Madau-Dickinson SFRD, log xi\_ion=25.5±0.15, clumping C=2-5, and JWST-SFRD tail. In comparison, indirect-proxy-inferred f\_esc from LzLCS O32/beta calibrations yields a value of 0.080 (+0.147/-0.051). The median delta between required and inferred values is -0.007 dex-frac (16-84\%: -0.150 to +0.067), with 46\% of systematic Monte Carlo simulations showing a shortfall.
\begin{figure}\centering\includegraphics[width=\columnwidth]{result.png}\caption{The reionization ionizing-photon-budget shortfall for star-forming galaxies is not robust to systematics at z\~{}9 (optimistic systematic corner)}\end{figure}
\section{Caveats}
It is essential to acknowledge the limitations of our approach. Our results are based on an automated, single-selection, uncalibrated measurement, which may not fully capture the complexities of reionization. The reliance on published values and calibrations introduces potential biases and uncertainties. Furthermore, the lack of direct observational data from surveys like JWST or SDSS restricts the robustness of our findings. Additionally, the systematic uncertainties in xi\_ion x clumping x proxy-calibration may dominate over statistical errors, highlighting the need for further research to refine these parameters.
\section*{References}
\footnotesize
[Muoz2024] Muñoz et al. (2024). Reionization after JWST: a photon budget crisis?. bibcode 2024MNRAS.535L..37M \par [Park2022] Park et al. (2022). Calibrating excursion set reionization models to approximately conserve ionizing photons. bibcode 2022MNRAS.517..192P \par [Duncan2015] Duncan et al. (2015). Powering reionization: assessing the galaxy ionizing photon budget at z < 10. bibcode 2015MNRAS.451.2030D \par [Davies2021] Davies et al. (2021). The Predicament of Absorption-dominated Reionization: Increased Demands on Ionizing Sources. bibcode 2021ApJ...918L..35D \par [Madau2017] Madau (2017). Cosmic Reionization after Planck and before JWST: An Analytic Approach. bibcode 2017ApJ...851...50M

\end{document}
